\scanpage{28}
\begingroup
\makeatletter\tagsleft@true\makeatother
à défaut d'appliquer le th. de dualité à $A$,
l'appliquer à $A'$, on trouve
\[
H^i(M)\simeq D\bigl(\operatorname{Ext}_{A'}^{\unsure{n'|N'}-i}(M,\Omega_{A'})\bigr)\qquad\text{i.e.}\tag*{(\deleted{\unsure{4}}1)}
\]
\[
\deleted{\operatorname{Ext}_{A'}^{\unsure{n'-i}}\unsure{k|(}}\tag*{\deleted{(1 bis)}}
\]
On en tire en particulier
\[
H^n(A)\simeq D\bigl(\operatorname{Ext}_{A'}^{n'-n}(M,\Omega_{A'})\bigr)\tag*{(2)}
\]
Ce qui prouve qu'on peut \deleted{\unsure{prendre}}\inserted{poser}
\[
\Omega_A\simeq\operatorname{Ext}_{A'}^{n'-n}(M,\Omega_{A'}).\tag*{(2 bis)}
\]

\deleted{Plus généralement, \unsure{introduisons}}

\[
\deleted{\Omega}\tag*{\deleted{(3)}}
\]
Soit $p$ un idéal premier de
$A$, \unsure{donc} $p\simeq A'/p'$, où $p'\supset\unsure{\mathfrak N|N}$. On a
alors
\[
(\Omega_A)_p\simeq\operatorname{Ext}_{A'_{p'}}^{n'-n}
(M_{p^{*}},(\Omega_{A'})_{p'}).
\]

Supposons qu'on sache que $(\Omega_{A'})_{p'}$ est
un module fond. pour l'anneau
[\uline{Coh. Mac.}] $A'_{p'}$ : ce sera le cas en
particulier si $A'$ est \uline{régulier} car alors
$\Omega_{A'}\simeq A'$, et $A'_{p'}$ sera régulier (Serre)
donc $A'_{p'}$ en sera un mod. fond. Alors
$(\Omega_A)_p$ est un mod. dualisant pour $\deleted{\unsure{A}}A_p$
d'après (2 bis) appliqué à $A_p\unsure{\Rightarrow|\simeq|=}A'_{p'}/\mathfrak N A'_{p'}$.
\endgroup
% archive sheet number: 27

